\documentclass[11pt]{article}
\usepackage{preamble}
\renewcommand{\answer}[1]{}

\begin{document}

\ladderheading{AP Statistics \textperiodcentered\ Unit 3 \textperiodcentered\ Topic 3.1 \textperiodcentered\ B: 2026-11-20 \textperiodcentered\ E: 2026-12-18}{Centered or Close?}

\focusbanner{TODAY'S PROBLEM}

Suppose a research team tests two rules for estimating a support proportion in a validation population where the true proportion is $p=0.40$. For a random sample of 25 people, let $X$ be the number who support the proposal.\par\smallskip Rule A reports $\widehat p_A=X/25$. Rule B reports $\widehat p_B=(X+5)/35$. Exact summaries of the two sampling distributions are shown.\par\smallskip \textbf{Lena:} ``A is unbiased, so it is always the more trustworthy estimator.''\par \textbf{Omar:} ``B varies less and is more likely to land within 0.06 of the truth, so its bias does not matter.''\par
\begin{center}
\textbf{Sampling distributions when $p=0.40$}\par\smallskip
\begin{tabular}{|l|c|c|c|}
\hline
\textbf{Rule} & \textbf{Mean} & \textbf{SD} & \textbf{Within 0.06} \\ \hline
\textbf{A} & 0.400 & 0.098 & 0.459 \\ \hline
\textbf{B} & 0.429 & 0.070 & 0.659 \\ \hline
\end{tabular}
\end{center}
\begin{firsttakebox}{FIRST TAKE --- before the video (5 min)}
Which rule looks more useful for getting close to 0.40? Use one entry in the table.
\workspace{0.9in}
\end{firsttakebox}

\begin{callout}{\IconBook\ BIAS AND VARIABILITY ANSWER DIFFERENT QUESTIONS (keep on the page)}{calloutgreen}
A \textbf{point estimator} is a statistic used to estimate a population parameter; its observed value is a
\textbf{point estimate}. An estimator is \textbf{unbiased} when the mean of its sampling distribution equals the
parameter. \textbf{Bias} is long-run center minus the parameter; \textbf{variability} is sampling-distribution spread.
Unbiased does not mean exact in one sample, and smaller spread does not erase bias. Judge the tradeoff against the stated goal.
\end{callout}

\newpage
\textbf{Finish after the video:}\par\smallskip

\dokbadge{1}~\textbf{(a)} If one sample has $X=9$ supporters, calculate the point estimate from Rule A and from Rule B. State the population parameter both rules estimate.\par
\workspace{0.8in}

\dokbadge{2}~\textbf{(b)} Use the table to classify each rule as biased or unbiased. For any biased rule, give the direction and amount of bias; then compare the rules' sampling variability.\par
\workspace{1.4in}

\dokbadge{3}~\textbf{(c)} In 3--4 sentences, choose a rule for getting within 0.06 of $p$ and use the table to defend it. Explain what Lena and Omar each overlook about bias and variability.\par
\workspace{2.0in}

\begin{sentenceframebox}\raggedright\textbf{Frame:} For the within-0.06 goal, I would use Rule \blank[0.4in] because \blank[2.6in].\end{sentenceframebox}

\begin{sentenceframebox}\raggedright\textbf{Frame:} That choice is not always better: its \blank[1.2in] could make a reader believe \blank[2.5in].\end{sentenceframebox}

\begin{summaryexitbox}{EXIT --- turn this in}
One thing the video changed about my first take: \hrulefill\par
\workspace{0.4in}
\end{summaryexitbox}

\end{document}
